\section{Incidencia ternaria, pegado de discriminantes y construcción de Leech}
\label{sec:lf-foundations}

La construcción reticular utiliza la información de incidencia que
acompaña al registro dodecafásico. Su dominio conserva las palabras
ternarias, el marco en que se representan, la orientación previa al
acarreo y la bandera que distingue una de las doce posiciones. El
paso a un retículo transmite esos datos mediante operaciones
sucesivas: codificación lineal, selección incidencial, pegado de
discriminantes y formación de un vecino. La clasificación reticular
interviene después de probar la positividad, paridad, autodualidad
y ausencia de raíces de la construcción obtenida.

Estas operaciones continúan la misma generación
APP--TRIT--TPK y el mismo refinamiento compatible del continuo
conjunto que anteceden a la lectura numérica \cite{HMT-IV}.
La matriz finita, el código y la métrica que se definen a
continuación son realizaciones con dominios precisos. La prueba
reticular recibe el estado marcado; el escalar final de una
constante no desempeña el papel de código de pegado ni de longitud.

\subsection{Coordenatización elíptica y transporte de la incidencia}
\label{subsec:lf-incidence}

Las seis unidades \(\{1,2,4,5,7,8\}\) módulo nueve, conservadas
por el transporte modular, admiten una coordenatización de seis posiciones.
La identificamos con
\(\mathbb P^1(\mathbb F_5)=\{\infty,0,1,2,3,4\}\), en ese orden.
Esta identificación fija coordenadas para realizar la relación
cuadrática del régimen elíptico del TRIT.

Sea \(\chi:\mathbb F_5\to\{0,1,-1\}\) el carácter cuadrático:
\(\chi(0)=0\), \(\chi(1)=\chi(4)=1\) y
\(\chi(2)=\chi(3)=-1\). Definimos la matriz entera
\(C_{\mathrm P}\) con diagonal cero,
\((C_{\mathrm P})_{\infty,a}=(C_{\mathrm P})_{a,\infty}=1\),
y entradas finitas \((C_{\mathrm P})_{a,b}=\chi(b-a)\).
Su cuadrado es \(5I_6\). En efecto, cada entrada diagonal del
cuadrado suma cinco unidades. Las entradas que involucran
\(\infty\) fuera de la diagonal se anulan porque la suma de
\(\chi\) es cero. Para dos índices finitos distintos, la
contribución uno de la posición \(\infty\) se cancela con
\(\sum_c\chi(c-a)\chi(c-b)=-1\). Esta última identidad se verifica
directamente sobre los cinco residuos, o mediante la sustitución
afín que lleva \((a,b)\) a \((0,1)\).

La reducción módulo tres produce la matriz simétrica
\begin{equation}
 A_W=
 \begin{pmatrix}
 0&1&1&1&1&1\\
 1&0&1&2&2&1\\
 1&1&0&1&2&2\\
 1&2&1&0&1&2\\
 1&2&2&1&0&1\\
 1&1&2&2&1&0
 \end{pmatrix},
 \qquad A_W^2=2I_6=-I_6
 \quad\text{en }\mathbb F_3.
 \label{eq:lf-aw}
\end{equation}
La igualdad conserva el tipo elíptico en este codominio finito.
El cuerpo finito y la posterior realización arquimediana de la
unidad imaginaria tienen dominios diferentes.

Para una matriz \(A\in M_6(\mathbb F_3)\) y una palabra fila
\(w\in\mathbb F_3^6\), escribimos
\[
 \gamma_A(w)=(w,wA)\in\mathbb F_3^{12}.
\]
Si \(f\in\operatorname{GL}_6(\mathbb F_3)\), entonces
\begin{equation}
 \gamma_{fAf^{-1}}(wf^{-1})
   =\gamma_A(w)(f^{-1}\oplus f^{-1}).
 \label{eq:lf-frame-covariance}
\end{equation}
La primera mitad es \(wf^{-1}\); la segunda es
\(wf^{-1}fAf^{-1}=wAf^{-1}\), lo que prueba la identidad.

Sea \(X_n^{\mathrm{mar}}\) el espacio de prefijos enriquecidos de
longitud \(n\), con marco inicial declarado, y sean
\(p_{n,m}:X_n^{\mathrm{mar}}\to X_m^{\mathrm{mar}}\) sus
truncaciones. En el paso \(j\), el prefijo contiene una palabra
visible \(w_j\) y un marco \(f_j\). La actualización del marco
tiene la forma \(f_{j+1}=g_j(x)f_j\), donde \(g_j(x)\) depende
únicamente de información ya contenida en el prefijo. Una
prolongación compatible conserva, por ello, todos los marcos
anteriores.

\begin{proposition}[Naturalidad del transporte incidencial]
\label{prop:lf-incidence-naturality}
La aplicación
\begin{equation}
 \mathcal Q_n(x)=
 \bigl((f_j,\gamma_{f_jA_Wf_j^{-1}}(w_j))\bigr)_{0\leq j<n}
 \label{eq:lf-Qn}
\end{equation}
satisface \(r_{n,m}\mathcal Q_n=\mathcal Q_mp_{n,m}\), donde \(r_{n,m}\) trunca
las primeras \(m\) componentes codificadas. Por tanto induce una
aplicación entre los límites inversos de los sistemas de
prefijos y de incidencias marcadas.
\end{proposition}
\begin{proof}
Dos prefijos compatibles tienen las mismas palabras y los mismos
operadores de transporte hasta el paso \(m-1\). Partiendo del
mismo marco inicial, la relación recursiva da por inducción los
mismos \(f_j\) hasta ese paso. Coinciden entonces las primeras
\(m\) componentes de \eqref{eq:lf-Qn}. La propiedad universal
del límite inverso produce la aplicación de prolongación.
\end{proof}

En cada marco, la primera proyección de \(\gamma_A\) recupera
exactamente la palabra visible. Este alcance es local: la
memoria profunda se conserva en el estado enriquecido y su
recuperación exige las coordenadas correspondientes. Al pasar
de palabras a soportes se retiene todavía menos información.
Asimismo, un cambio general de base en
\(\operatorname{GL}_6(\mathbb F_3)\) transporta también la métrica
de Hamming y la incidencia de la coordenatización. Los cambios monomiales
las preservan directamente. Los soportes utilizados en esta
sección se calculan en la coordenatización explícita \eqref{eq:lf-aw}.

La naturalidad anterior concreta la lectura incidencial que
acompaña a los cilindros numéricos. El encajamiento de cilindros
y la restricción de \eqref{eq:lf-Qn} actúan sobre un mismo
prefijo compatible, conservando respectivamente su evaluación
posicional y sus relaciones marcadas. El ambiente
\(\mathbb F_3^6\), de \(729\) palabras, proporciona el dominio
lineal de la codificación; conserva su distinción respecto de
las \(243\) palabras visibles y las \(468\) emisiones del censo
admisible de APP--TRIT--TPK.

\subsection{Código ternario autodual y enumeración de sus pesos}
\label{subsec:lf-code}

\begin{definition}[Código de la coordenatización dodecafásica]
El código lineal asociado a \eqref{eq:lf-aw} es
\[
 C_W=\{(w,wA_W):w\in\mathbb F_3^6\}\subset\mathbb F_3^{12}.
\]
El peso \(\operatorname{wt}(c)\) cuenta las coordenadas no nulas
de \(c\), y \(\operatorname{supp}(c)\subset\{1,\ldots,12\}\)
es su conjunto de posiciones no nulas. Una hexada es el soporte
de una palabra de peso seis; su colección se denota \(\mathcal H\).
\end{definition}

\begin{theorem}[Autodualidad, distancia mínima y enumerador]
\label{thm:lf-code}
El código \(C_W\) es autodual de parámetros \([12,6,6]_3\) y
tiene enumerador
\begin{equation}
 W_{C_W}(z)=1+264z^6+440z^9+24z^{12}.
 \label{eq:lf-enumerator}
\end{equation}
\end{theorem}
\begin{proof}
La matriz generadora \([I_6\mid A_W]\) tiene rango seis.
Por la simetría de \(A_W\) y \eqref{eq:lf-aw},
\[
 [I_6\mid A_W][I_6\mid A_W]^{\mathsf T}
       =I_6+A_W^2=0.
\]
Así, el código es autoortogonal de dimensión mitad de la
longitud y, en consecuencia, autodual.

La enumeración restante es una operación entera finita sobre
la matriz explícita. Para \(r\in\{0,\ldots,6\}\), agrupamos
las palabras por \(\operatorname{wt}(w)=r\). La tabla siguiente
da el número de palabras con cada peso total:
\[
\begin{array}{c|rrrr|r}
r&0&6&9&12&\text{total}\\ \hline
0&1&0&0&0&1\\
1&0&12&0&0&12\\
2&0&60&0&0&60\\
3&0&120&40&0&160\\
4&0&60&180&0&240\\
5&0&12&180&0&192\\
6&0&0&40&24&64\\ \hline
\text{total}&1&264&440&24&729
\end{array}
\]
Cada fila contiene \(\binom6r2^r\) palabras. Para verificar
todas sus entradas, sea \(q\) un entero de \(0\) a \(728\), y
definamos sus seis cifras ternarias
\[
 d_i(q)=\left\lfloor\frac{q}{3^{6-i}}\right\rfloor\bmod3,
 \qquad
 e_j(q)=\left(\sum_{i=1}^{6}d_i(q)(A_W)_{ij}\right)\bmod3.
\]
La fila y la columna de la contribución de \(q\) son,
respectivamente,
\[
 r(q)=\sum_i\mathbf1_{\{d_i(q)\ne0\}},\qquad
 h(q)=r(q)+\sum_j\mathbf1_{\{e_j(q)\ne0\}}.
\]
La recurrencia de acumulación empieza con
\(T^{(0)}_{r,h}=0\) y aplica
\[
 T^{(q+1)}_{r,h}
 =T^{(q)}_{r,h}
   +\mathbf1_{\{r=r(q)\}}\mathbf1_{\{h=h(q)\}}
 \quad(0\leq q<729).
\]
Produce exactamente la tabla mostrada. La expansión ternaria
da una biyección entre esos \(729\) enteros y todas las palabras
de \(\mathbb F_3^6\), por lo que la operación agota el código.
La suma de sus filas prueba \eqref{eq:lf-enumerator}; el primer
peso no nulo es seis.
\end{proof}

\begin{remark}[Evaluación reproducible de la recurrencia finita]
El siguiente bucle implementa literalmente las sumas de la prueba.
Sus únicos datos son los enteros de la matriz
\eqref{eq:lf-aw}; calcula el enumerador antes de compararlo con
la igualdad obtenida.
\end{remark}
{\small
\begin{verbatim}
from collections import Counter
from itertools import product
A = ((0,1,1,1,1,1), (1,0,1,2,2,1),
     (1,1,0,1,2,2), (1,2,1,0,1,2),
     (1,2,2,1,0,1), (1,1,2,2,1,0))
rows = [Counter() for _ in range(7)]
for w in product(range(3), repeat=6):
    v = tuple(sum(w[i]*A[i][j] for i in range(6)) % 3
              for j in range(6))
    r = sum(x != 0 for x in w)
    h = r + sum(x != 0 for x in v)
    rows[r][h] += 1
total = sum(rows, Counter())
if dict(total) != {0:1, 6:264, 9:440, 12:24}:
    raise ArithmeticError("Enumerador incompatible")
\end{verbatim}
}

\begin{proposition}[Hexadas y diseño de Witt]
\label{prop:lf-witt}
La colección \(\mathcal H\) contiene \(132\) hexadas y cada
quintupla de posiciones pertenece a una única hexada.
\end{proposition}
\begin{proof}
Dos palabras de peso seis con el mismo soporte difieren por
un escalar. En efecto, entre sus seis cocientes de coordenadas
no nulas, uno de los signos \(1,-1\) aparece al menos tres
veces. Restando el múltiplo correspondiente, si las palabras
no fueran proporcionales, se obtendría una palabra no nula
de peso a lo sumo tres, en contradicción con
el teorema~\ref{thm:lf-code}. Cada soporte corresponde así a
la pareja \(\{c,-c\}\), y hay \(264/2=132\) soportes.

Si dos soportes distintos compartieran cinco posiciones,
en esa intersección uno de los dos cocientes aparecería al
menos tres veces. La combinación lineal que cancela dichas
coordenadas tendría soporte en una unión de tamaño siete
y peso a lo sumo cuatro. La distancia mínima lo impide.
Una quintupla pertenece, por tanto, a lo sumo a una hexada.
Como
\[
 132\binom65=792=\binom{12}{5},
\]
pertenece a exactamente una.
\end{proof}

El objeto construido es el código ternario extendido de Golay
con su pequeño diseño de Witt. El reconocimiento se efectúa
después de la codificación, la autodualidad y el cómputo de
distancia. Para la construcción reticular sólo se necesitan
las propiedades ya demostradas y la bandera marcada que se
obtiene a continuación.

\subsection{Bandera dodecafásica y selección del origen}
\label{subsec:lf-flag}

En el marco del registro generado se conservan las palabras
\[
 w_\pi=010211,\qquad
 w_e=201101,\qquad
 w_\varphi=121200,\qquad
 w_{\mathrm{APP}}=022020.
\]
Los subíndices \(\pi,e,\varphi\) identifican las correspondientes
regiones de lectura de las constantes; \(w_e\) conserva aquí
la región de Euler. Aplicar \(\gamma_{A_W}\) y tomar soportes da
\[
\begin{aligned}
 P_\pi&=\{2,4,5,6,7,8,9,10,11\},\\
 H_e&=\{1,3,4,6,9,10\},\\
 H_\varphi&=\{1,2,3,4,7,9\},\\
 H_{\mathrm{APP}}&=\{2,3,5,10,11,12\}.
\end{aligned}
\]
El primer soporte tiene peso nueve; los tres restantes son
hexadas. El registro
\[
 K=(234,543,140,729,659,824,621,58,914,794,146,601)
\]
selecciona
\[
 B_K=\{j:K_j\geq729\}=\{4,6,9,10\}=H_e\cap P_\pi.
\]
El umbral \(729\) pertenece al formato de bloques de la
lectura declarada. En este registro concreto, los umbrales
enteros de \(660\) a \(729\) producen el mismo soporte, de
modo que el soporte conserva la incidencia seleccionada,
mientras el registro completo conserva también los bloques
y el umbral de su lectura.

El registro firmado anterior al acarreo es
\[
\begin{split}
 Z_0={}&(7,297,353,-430,-715,-1199,\\
       &\hspace{9mm}-3,286,-894,-528,380,663).
\end{split}
\]
Su soporte negativo es la hexada
\[
 H_\alpha^-=\{4,5,6,7,9,10\}.
\]
Los signos se reciben del estado previo al acarreo y no se
extraen de la expansión del escalar final.

\begin{proposition}[Coincidencia de selección firmada e incidencial]
\label{prop:lf-flag}
La hexada \(H_\alpha^-\) es la única \(H\in\mathcal H\) que cumple
\[
\begin{gathered}
 B_K\subset H\subset P_\pi,\\
 (|H\cap H_e|,|H\cap H_\varphi|,
                   |H\cap H_{\mathrm{APP}}|)=(4,3,2).
\end{gathered}
\]
\end{proposition}
\begin{proof}
La codificación de las palabras de peso seis, o la recurrencia
finita anterior aplicada a sus soportes, da las cuatro
hexadas sobre \(B_K\). Sus pares complementarios son
\[
 \{1,3\},\qquad\{2,11\},\qquad\{5,7\},\qquad\{8,12\}.
\]
Sólo el segundo y el tercero están contenidos en \(P_\pi\).
La hexada del segundo tiene intersecciones de tamaño tres
con \(H_\varphi\) y con \(H_{\mathrm{APP}}\), y por tanto
incumple la última condición. La del tercero tiene
intersecciones \(4,3,2\), y coincide con \(H_\alpha^-\).
\end{proof}

\begin{lemma}[Origen de la bandera marcada]
\label{lem:lf-origin}
El origen seleccionado por la incidencia es
\[
 o_*=(H_e\cap H_\varphi)
                 \setminus(H_\alpha^-\cup H_{\mathrm{APP}})
     =\{1\}.
\]
La selección conmuta con toda biyección aplicada
simultáneamente a los soportes.
\end{lemma}
\begin{proof}
La intersección \(H_e\cap H_\varphi\) es
\(\{1,3,4,9\}\). La unión sustraída contiene \(3,4,9\)
y excluye \(1\). La equivariancia procede de que las
biyecciones conmutan con unión, intersección y diferencia.
\end{proof}

La bandera distingue una componente entre doce antes del
paso métrico. El vector radial se determina después,
dentro de una cámara positiva declarada. Así quedan
separadas y articuladas la selección combinatoria y la
operación reticular.

\subsection{Pegado ternario de las doce componentes}
\label{subsec:lf-glue}

Sea \(A_2=\mathbb Za_1\oplus\mathbb Za_2\), con forma
bilineal de matriz
\[
 B_2=\begin{pmatrix}2&-1\\-1&2\end{pmatrix},
 \qquad \rho=a_1+a_2,\qquad \rho^2=2.
\]
Para \(x=ua_1+va_2\),
\[
 x^2=2(u^2-uv+v^2).
\]
La identidad
\(4(u^2-uv+v^2)=(2u-v)^2+3v^2\) muestra
la positividad. El determinante es tres. Las tres
clases de \(A_2^*/A_2\) se representan mediante
\[
 \varpi_0=0,\qquad
 \varpi_1=\frac{2a_1+a_2}{3},\qquad
 \varpi_2=\frac{a_1+2a_2}{3}.
\]
El emparejamiento con \(a_1,a_2\) verifica que estos
vectores pertenecen a \(A_2^*\), y
\(2\varpi_1-\varpi_2=a_1\). Por ello
\(j\mapsto\varpi_j+A_2\) identifica aditivamente
\(\mathbb F_3\) con el discriminante.

Para \(c\in C_W\), definimos
\(\phi(c)=(\varpi_{c_1},\ldots,\varpi_{c_{12}})\)
y
\begin{equation}
 N=\bigcup_{c\in C_W}\bigl(A_2^{12}+\phi(c)\bigr).
 \label{eq:lf-glue}
\end{equation}

\begin{theorem}[Retículo de pegado par y unimodular]
\label{thm:lf-niemeier}
El conjunto \(N\) es un retículo positivo, par y
unimodular de rango \(24\). Sus vectores de norma dos
son exactamente las raíces de \(A_2^{12}\).
\end{theorem}
\begin{proof}
La linealidad del código y la aditividad del discriminante
hacen que la unión sea un subgrupo de
\((A_2^*)^{12}\) que contiene \(A_2^{12}\).
Es, por tanto, un retículo de rango \(24\).

El emparejamiento de las clases discriminantes es
\(2cd/3\) módulo \(\mathbb Z\) en cada componente.
La autoortogonalidad de \(C_W\) anula su suma y prueba
la integralidad de \(N\). La norma de una clase
representada por \(\phi(c)\) es
\(2\operatorname{wt}(c)/3\) módulo \(2\mathbb Z\).
Todos los pesos de \eqref{eq:lf-enumerator} son
múltiplos de tres, luego \(N\) es par.

Su índice sobre \(A_2^{12}\) es \(|C_W|=3^6\).
Con el determinante de la matriz de Gram,
\[
 \det N=\frac{\det(A_2^{12})}{[N:A_2^{12}]^2}
       =\frac{3^{12}}{(3^6)^2}=1.
\]
La forma ambiente es la suma ortogonal de doce formas
positivas; queda demostrada la positividad.
La integralidad y el determinante uno implican
\(N^*=N\).

En una clase no nula del discriminante local, un
vector tiene la forma \((ua_1+va_2)/3\), con residuos
\((u,v)\equiv(2,1)\) o \((1,2)\pmod3\).
El entero \(u^2-uv+v^2\) es positivo y divisible
por tres, luego es al menos tres. Los representantes
\(\varpi_1,\varpi_2\) alcanzan ese valor: la norma
mínima de cada clase no nula es \(2/3\).
Una palabra no nula del código tiene por lo menos
seis coordenadas no nulas. Todo vector de una clase
de pegado no trivial tiene, por tanto, norma
al menos \(6(2/3)=4\).

En \(A_2^{12}\), una norma dos sólo puede proceder
de una única componente, porque cada componente
tiene norma par positiva mínima dos. Las soluciones
de \(u^2-uv+v^2=1\) son
\((\pm1,0),(0,\pm1),(1,1),(-1,-1)\).
Son las seis raíces de \(A_2\), y completan el censo
de raíces de \(N\).
\end{proof}

Las propiedades anteriores identifican \(N\) como el
retículo de Niemeier de sistema de raíces \(A_2^{12}\).
La descripción de todas sus raíces es necesaria para
controlar cuáles sobreviven al paso al vecino.

\subsection{Cámara radial y vecino sin raíces}
\label{subsec:lf-neighbour}

Se fija la cámara radial
\[
 v(a)=\sum_{j=1}^{12}a_j\rho_j,\qquad
 a_j=1+3b_j,\qquad b_j\in\mathbb Z_{\geq0},
\]
donde \(\rho_j\) es la copia de \(\rho\) en la
componente \(j\). La paridad de un vecino de índice
tres requiere \(v(a)^2\equiv0\pmod{18}\).
Como
\[
 v(a)^2=2\sum_j(1+3b_j)^2
       =24+12\sum_jb_j+18\sum_jb_j^2,
\]
la congruencia equivale a \(\sum_jb_j\equiv1\pmod3\).
La norma mínima dentro de la cámara se alcanza con
un único \(b_j=1\) y los demás cero, y vale \(54\).
El origen \(o_*=1\) del lema~\ref{lem:lf-origin}
selecciona
\begin{equation}
 v_\alpha=4\rho_1+\rho_2+\cdots+\rho_{12},
 \qquad v_\alpha^2=54.
 \label{eq:lf-valpha}
\end{equation}
El calificativo de mínimo se refiere a la cámara
declarada. En la clase residual completa,
\((a_1-2a_2,\rho,\ldots,\rho)\) representa la
misma clase módulo tres y tiene norma
\(14+11\cdot2=36\). Esta distinción preserva
la condición que determina el coeficiente cuatro.

Definimos
\begin{equation}
\begin{split}
 N_0&=\{x\in N:\langle x,v_\alpha\rangle\equiv0\pmod3\},\\
 \Lambda&=N_0+\mathbb Z\,v_\alpha/3.
\end{split}
\label{eq:lf-lambda}
\end{equation}

\begin{theorem}[Vecino par autodual sin raíces]
\label{thm:lf-leech}
La construcción \eqref{eq:lf-lambda} es un
retículo positivo, par y autodual de rango \(24\)
sin vectores de norma dos. En consecuencia, es
isométrica al retículo de Leech.
\end{theorem}
\begin{proof}
El carácter \(x\mapsto\langle x,v_\alpha\rangle\bmod3\)
es no nulo. Una raíz simple de la componente \(j\)
tiene emparejamiento \(a_j\equiv1\pmod3\) con
\(a_j\rho_j\). Por tanto \([N:N_0]=3\), de donde
\(\det N_0=9\).

Para \(x\in N_0\), el emparejamiento
\(\langle x,v_\alpha/3\rangle\) es entero, y así
\(v_\alpha/3\in N_0^*\). Además,
\(\langle v_\alpha,v_\alpha\rangle=54\) implica
\(v_\alpha\in N_0\). Si \(v_\alpha/3\in N\),
la integralidad de \(N\) obligaría a que todos los
emparejamientos de \(N\) con \(v_\alpha\) fuesen
divisibles por tres, contradiciendo el carácter
no nulo. La clase de \(v_\alpha/3\) sobre \(N_0\)
tiene entonces orden exactamente tres. Resulta
\[
 [\Lambda:N_0]=3,\qquad
 \det\Lambda=\frac9{3^2}=1.
\]
Los emparejamientos de \(N_0\) con el nuevo
generador son enteros y la norma de este generador
es seis. La extensión es integral. Para
\(x\in N_0\) y \(m\in\mathbb Z\),
\[
 \left(x+m\frac{v_\alpha}{3}\right)^2
 =x^2+2m\frac{\langle x,v_\alpha\rangle}{3}
       +6m^2\in2\mathbb Z.
\]
Por ello \(\Lambda\) es par; su integralidad y
determinante uno prueban la autodualidad.
La positividad y el rango se heredan del espacio
euclídeo ambiente.

Para excluir raíces, primero consideramos la
clase \(N_0\). Todas las raíces de \(N\) pertenecen
a \(A_2^{12}\) por el teorema~\ref{thm:lf-niemeier}.
Los emparejamientos de sus seis raíces locales
con \(a_j\rho_j\) son \(\pm a_j,\pm2a_j\),
congruentes con \(\pm1,\pm2\pmod3\).
Ninguna raíz pertenece a \(N_0\).

Las otras dos clases de \(\Lambda/N_0\) pueden
representarse con \(\pm v_\alpha/3\).
En cada componente, sus vectores pertenecen a
una de las clases
\[
 A_2+\varpi_j\pm\rho/3,\qquad j=0,1,2.
\]
La componente distinguida tiene la misma
descripción porque \(4\rho/3\equiv\rho/3\pmod{A_2}\).
Para los numeradores \((u,v)\) en
\((ua_1+va_2)/3\), los seis residuos son
\[
 (1,1),(0,2),(2,0),(2,2),(1,0),(0,1)\pmod3.
\]
Todos son no nulos. La positividad de
\(u^2-uv+v^2\) obliga a que sea al menos uno,
y los representantes
\((1,1),(0,-1),(-1,0),(-1,-1),(1,0),(0,1)\)
alcanzan ese mínimo en las seis clases.
Cada componente tiene norma al menos \(2/9\).
Un vector de cualquiera de las dos clases
nuevas tiene norma al menos
\[
 12\cdot\frac29=\frac83>2.
\]
Se han excluido raíces en las tres clases.
El teorema de clasificación de los retículos
positivos pares unimodulares de rango veinticuatro
identifica el único caso sin raíces con el
retículo de Leech \cite{ConwaySloane}.
\end{proof}

La identificación final utiliza una clasificación
externa sobre un objeto cuyas propiedades ya han
sido demostradas. La selección de la componente
distinguida conserva la bandera procedente del
registro firmado; el subíndice \(\alpha\) expresa
esa procedencia. La norma y los coeficientes del
vector radial se obtienen de la métrica y de
la cámara mediante la congruencia indicada.

El retículo \(\Lambda\), su forma bilineal, sus
doce planos ambientes y el vector \(v_\alpha/3\)
quedan así determinados con todas las propiedades
utilizadas por la deformación posterior. El
código contiene asimismo la palabra
\[
 c_*=(1,1,1,1,1,1,2,1,1,1,1,1),
\]
pues \(q_*=(1,1,1,1,1,1)\) satisface
\(q_*A_W=(2,1,1,1,1,1)\). Éste es el dato
de soporte completo que permite construir y
probar la estabilidad de la isometría de orden
tres en la sección~\ref{sec:leech-dualidad}.
